The distribution is organized into three top-level folders.

\subsection{Geographic taxonomy files (\texttt{geo/})}

Two files provide the geographic taxonomy of local units used in all results files.

\subsubsection{\texttt{geo/laus.csv}}

Lists municipality-level units defined by Eurostat's
Geographic Information System of the Commission (GISCO).
Each entry in this file corresponds to one Local Administrative Unit (LAU),
the most fine-grained level of European
administrative units. BLUE\_DB tracks vertical and horizontal dependencies between
these units, recording both parent regions from the
Nomenclature of Territorial Statistical Units (NUTS)
and temporal predecessor/successor relationships resulting from municipality
mergers and splits.

\subsubsection{\texttt{geo/nuts.csv}}

Lists regions-level units defined by Eurostat's GISCO.
Each entry in this file corresponds to one NUTS region.
As with municipalities, BLUE\_DB similarly tracks vertical and horizontal dependencies
between regional units.

\subsubsection{\texttt{geo/special.csv}}

Contains special-purpose units that are needed to report election results
but do not correspond to a standard LAU or NUTS region. Examples include postal-vote
units covering multiple municipalities, overseas voter counts,
various electoral constituencies above the municipal tier, and regional
aggregate rows. Each row carries a \texttt{type} field.

\subsection{Party taxonomy (\texttt{parties/})}

\subsubsection{\texttt{parties/parties.csv}}

Lists all parties, lists, coalitions, and independent candidates referenced across the dataset,
providing its official name(s), dates of creation and dissolution, European
affiliations to a European party or EP group, broad ideological position,
and various identifiers.

\subsubsection{\texttt{parties/eu\_parties.csv}}

Lists all European parties referenced across the dataset, providing similar information
as for national or regional parties.

\subsubsection{\texttt{parties/ep\_groups.csv}}

Lists all European Parliament groups references across the dataset, providing similar
information as for national or regional parties.

\subsection{Election files (\texttt{elections/})}

This folder contains one subfolder by country covered, named by its country
code. Each country folder contains a file
\texttt{YEAR\_MONTH.csv} for national elections and \texttt{EU\_YEAR\_MONTH.csv} for
European Parliament elections, with suffixes where multiple ballots, rounds of voting,
or vote counts exist in one year.
The special file \texttt{elections/index.csv} lists all available results files.

Party columns are \textit{wide}: each distinct \texttt{party\_id} that
appears in that election file occupies one column; cells are empty where the
party did not stand.
All vote figures are \textit{absolute vote counts} as published in official
sources unless noted otherwise in the country-specific section.
